% =============================================================================
% SPRS Hardware Implementation Reference — arXiv companion to the TPDS paper
% "Bitwise-Reproducible Reduction Across Cluster Reshape: A Compiler/Fabric ABI"
%
% Source material: design_choices.md (module-by-module rationale) and the RTL
% itself. This document carries the hardware detail the main paper cites but
% has no page budget for.
% =============================================================================
\documentclass[journal,10pt]{IEEEtran}

\usepackage{amsmath,amssymb,amsfonts}
\usepackage{booktabs}
\usepackage{graphicx}
\usepackage{xcolor}
\usepackage{listings}
\usepackage{url}
\usepackage[hidelinks]{hyperref}

\newcommand{\sprs}{\textsc{sprs}}
\newcommand{\sv}[1]{\texttt{#1}}
\newcommand{\FBT}{\mathrm{FBT}}

\lstset{basicstyle=\ttfamily\scriptsize, breaklines=true,
        columns=fullflexible, frame=single, framesep=3pt,
        rulecolor=\color{gray!50}, showstringspaces=false}

\begin{document}

\title{SPRS: Hardware Implementation Reference\\
{\large A companion to ``Bitwise-Reproducible Reduction Across Cluster
Reshape: A Compiler/Fabric ABI''}}

\author{%
Siddharth~Patel,
Aman~Awasthi,
Akshansh~Chaurasia,
and~Rohit~Singh%
\thanks{The authors are with the Department of Electrical Engineering,
Shiv Nadar University, India.
E-mail: \{sp839, aa634, ac198, rohit.singh\}@snu.edu.in}%
}

\maketitle

\begin{abstract}
This companion document specifies the \sprs{} hardware in the detail a
reimplementation would require: the register-transfer structure of every
module, the complete instruction and packet encodings, the routing-table
format, the arbitration and flow-control disciplines, the FPGA bring-up
procedure, and the verification strategy. It is a reference, not an argument.
The claims about bitwise reproducibility, the ABI that produces them, and
their experimental validation are made in the main paper; here we document
only what the fabric \emph{is}. Readers wanting the result should start with
the main paper and return here for implementation specifics.
\end{abstract}

\begin{IEEEkeywords}
Network-on-chip, reduction tree, deterministic arithmetic, IEEE 754,
SystemVerilog, FPGA.
\end{IEEEkeywords}

% =====================================================================
\section{Scope and Organisation}
% =====================================================================

The \sprs{} fabric is 3{,}726 lines of SystemVerilog across thirteen
files: ten implementing the fabric proper and three FPGA wrappers. This
document walks them in dependency order --- packages, primitives, link layer,
router, network interface, compute node, top level --- and then covers the
cross-cutting decisions that do not belong to any single module.

\begin{table}[!t]
\caption{Source inventory.}
\label{tab:inventory}
\centering
\footnotesize
\begin{tabular}{@{}llr@{}}
\toprule
File & Role & Lines \\
\midrule
\sv{noc\_pkg.sv}            & NoC types, packet layout      & 41 \\
\sv{btree\_pkg.sv}          & ISA, ALU ops, parameters      & 101 \\
\sv{sync\_fifo.sv}          & FIFO primitive                & 61 \\
\sv{link\_tx.sv}            & credit-based transmit         & 37 \\
\sv{noc\_link.sv}           & inter-router link, latency    & 78 \\
\sv{noc\_router.sv}         & $N$-port multi-VC router      & 295 \\
\sv{noc\_ni.sv}             & network interface, GCI buffer & 217 \\
\sv{fp64\_add.sv}           & IEEE-754 FP64 adder           & 242 \\
\sv{btree\_fsm\_fast.sv}    & per-PE reduction pipeline     & 912 \\
\sv{noc\_system.sv}         & topology-agnostic top level   & 303 \\
\midrule
\sv{top\_fpga\_interactive.sv}     & Artix-7 wrapper, $G{=}2$   & 503 \\
\sv{top\_fpga\_interactive\_g2.sv} & deterministic $G{=}2$      & 430 \\
\sv{top\_fpga\_interactive\_g4.sv} & deterministic $G{=}4$      & 506 \\
\midrule
\multicolumn{2}{@{}l}{\textbf{Total}} & \textbf{3{,}726} \\
\bottomrule
\end{tabular}
\end{table}

% =====================================================================
\section{Instruction and Packet Encodings}
% =====================================================================

\subsection{Instruction word}

Instructions are 64 bits, four opcodes, fixed fields:

\begin{lstlisting}
[63:60] opcode   NOP=0 GENERATE=1 COMPUTE=2 SEND=3
[59:56] aux      aux[0]=link, aux[3:1]=alu_op
[55:40] dest     DMEM destination (COMPUTE/GENERATE)
                 target_addr at receiver (SEND)
[39:24] addr_a   DMEM source A (COMPUTE)
                 local DMEM source (SEND)
[23:8]  addr_b   DMEM source B (COMPUTE)
                 dest_gpu id (SEND)
[7:0]   reserved
\end{lstlisting}

Opcode~4 is \sv{OP\_RESERVED}, formerly \sv{OP\_RECV}, retired when the
messaging model moved to RDMA-push. Its retirement is load-bearing rather
than cosmetic: with no receive instruction there is no receive-side blocking
construct, and therefore no way for a program to deadlock waiting on a
message that never arrives. Delivery is the network's obligation, not the
program's.

The ALU field admits eight operations (\sv{ADD}, \sv{MAX}, \sv{MIN},
\sv{AND}, \sv{OR}, \sv{XOR}, \sv{SADD}, \sv{FADD}). Only \sv{FADD}
(IEEE-754 double addition) is exercised in the main paper's experiments; the
others are implemented and functional but unvalidated at scale, and should be
treated as untested surface.

\subsection{Packet format}

Packets are single-flit, 96 bits:

\begin{lstlisting}
[95:80] dest_node_id  (16b) target PE, used for routing
[79:70] target_addr   (10b) DMEM address at the destination
[69:64] reserved      (6b)
[63:0]  payload       (64b) the value
\end{lstlisting}

Single-flit packets are the reason the fabric needs no reassembly logic,
no flit sequencing, and no virtual-cut-through buffering, at the cost of a
96-bit datapath width. Bit~63 of the payload additionally steers the virtual
channel, a deliberate choice discussed in \S\ref{sec:vc}.

\subsection{Routing tables}

Routing is table-driven: router $r$ holds \sv{route\_table[d]} giving the
egress port for destination $d$. Tables are computed by the compiler and
loaded at configuration time. This makes topology a configuration input
rather than a structural property of the RTL --- the same netlist serves
mesh, torus, hypercube, fat-tree, ring and linear arrangements, selected by
the adjacency table passed to \sv{noc\_system}.

A width subtlety costs correctness if missed. \sv{route\_table} elements are
$\lceil\log_2(\text{NUM\_PORTS})\rceil$ bits wide. Most topologies use
five-port routers, giving three bits, but fat-tree routers have more ports
and need four. Any tooling that writes the table with a sized literal
(\sv{3'd}$p$) silently truncates port numbers above seven and breaks fat-tree
routing in a way that passes elaboration and fails only at run time. Plain
integer literals, which widen to the left-hand side, are required.

% =====================================================================
\section{Compute Node: \texttt{btree\_fsm\_fast}}
% =====================================================================

The per-PE reduction engine is a three-stage pipeline sustaining one
instruction per cycle.

\begin{center}
\begin{tabular}{@{}lll@{}}
\toprule
Stage & Reads & Does \\
\midrule
IF & IMEM        & fetch \\
DE & DMEM        & operand read, hazard check \\
EX & ALU / TX    & compute, enqueue send, write back \\
\bottomrule
\end{tabular}
\end{center}

\subsection{Scoreboard gating}

Every DMEM address carries a presence bit. \sv{DE} stalls if an instruction
reads an address whose bit is clear. This is the mechanism that makes
arrival order unobservable to the program: a \sv{COMPUTE} reads its operands
\emph{by address}, and proceeds only once both have been written, regardless
of which agent wrote them, in what order, over which virtual channel, or at
what time. The scoreboard is what converts a network with nondeterministic
timing into a machine with deterministic values.

The presence bit is monotone within a program run --- set on write, never
cleared --- so the gate cannot oscillate.

\subsection{Write arbitration}

Three agents write DMEM: the packet ingress unit (PIU), the GCI leaf
interface, and the FSM writeback. Priority is fixed at
$\text{PIU} > \text{GCI} > \text{FSM}$. PIU wins because a dropped inbound
packet is unrecoverable --- the sender has already retired the instruction ---
whereas the FSM can stall a cycle without losing information. The GCI path
carries a write-hold latch so that a leaf value colliding with a PIU write is
retained rather than discarded.

\subsection{Forwarding}

One level of write-to-read forwarding covers back-to-back \sv{COMPUTE}
instructions, where the second consumes the first's result. Without it the
scoreboard would stall every dependent pair for a cycle. The forwarding path
is value-carrying only; it does not bypass the presence bit, so it cannot
manufacture a read of never-written data.

\subsection{Fail-stop discipline}

Three independent mechanisms convert malformed state into a halt rather than
a wrong answer: a per-PE stall timeout, a global watchdog, and an
invalid-opcode trap. A misrouted packet arriving at the wrong PE is dropped
rather than written, so a routing-table error manifests as a watchdog
timeout, never as a silently corrupted result. The main paper states this as
Property~2.

% =====================================================================
\section{Router}\label{sec:vc}
% =====================================================================

An $N$-port router with \sv{NUM\_VCS} virtual channels per input, default
two, in a three-stage arrangement: input buffering per port per VC, route
computation with VC and switch allocation, then switch traversal. Arbitration
is two-level round-robin --- among VCs within a port, then among ports for
each output --- and advances its pointer only on a successful grant, so a
blocked request cannot starve its neighbours by consuming arbitration turns.

Virtual channels exist for deadlock freedom on cyclic topologies. The VC is
selected by packet bit~63, that is, by the compiler rather than by hardware
policy. This keeps the router simple and puts deadlock avoidance where the
global view lives, but it also means a compiler bug can construct a cyclic
channel dependency the hardware will faithfully deadlock on. The watchdog
converts that into a fail-stop.

Flow control is credit-based, with an eight-bit credit counter per link.
Credit returns are counter-serialised: an earlier revision pulsed credits
combinationally and could lose a return when two arrived in the same cycle,
leaking credits until the link wedged.

% =====================================================================
\section{Arithmetic}
% =====================================================================

\sv{fp64\_add} is a fully combinational IEEE-754 double-precision adder:
round-to-nearest-even, flush-to-zero on both subnormal inputs and subnormal
results. It handles normal$+$normal, signed zeros, infinities including
$\infty + (-\infty) \rightarrow$ NaN, and NaN propagation.

It is stateless and unpipelined, and both properties are deliberate. A
stateless adder is a pure function of its two operands, so the same operand
pair yields the same 64-bit result regardless of what the unit computed
previously --- no accumulated rounding mode, no sticky flags, no history. An
adder carrying mode state would break the main paper's Claim~1 at the
arithmetic level no matter what the compiler and network did correctly.

Flush-to-zero is a genuine deviation from full IEEE-754 conformance and is
declared as such. It is uniform across every configuration, so it does not
threaten reproducibility, but it means \sprs{} results need not agree with an
external reducer that honours subnormals.

% =====================================================================
\section{Resource Envelope and Known Limits}
% =====================================================================

Fixed capacities: DMEM and IMEM 512 entries per PE, TX and RX buffers 8
entries, timeout 10{,}000 cycles, watchdog 20{,}000. \sv{MAX\_ROUTERS} is
4096.

The 512-entry instruction memory is the binding constraint on instance
feasibility, and it is the direct cause of the IMEM-overflow compilation
failures reported in the main paper: a partition that assigns too many tree
nodes to one PE produces an instruction image that does not fit, and the
compiler rejects it rather than emitting a truncated program.

Single-flit packets combined with input-buffered routers mean head-of-line
blocking is possible at large $G$ on low-bisection topologies. Wormhole
forwarding would address it and is not implemented.

% =====================================================================
\section{FPGA Bring-Up}
% =====================================================================

Three wrappers target a Xilinx Artix-7 (XC7A100T, Nexys-family board). The
interactive wrapper combines a hardware bootloader that loads instruction
memory, routing tables and adjacency at reset; VIO probes for entering leaf
values and triggering a reduction; an ILA for observing pipeline activity;
and a seven-segment display of the resulting root.

Procedure:
\begin{enumerate}
\item Program the bitstream; the config FSM loads IMEM, routing and adjacency.
\item Drive \sv{vio\_leaf\_value} and \sv{vio\_leaf\_idx}, pulse
      \sv{vio\_leaf\_wr}, once per leaf.
\item Pulse \sv{vio\_start} or flip \sv{SW[0]}.
\item Read \sv{vio\_result} (64-bit root) and \sv{vio\_perf\_cycles}.
\end{enumerate}

When recording cycle counts, take the maximum of \sv{perf\_total} across all
PEs, not PE~0's counter. The root of the reduction frequently lands on a PE
other than zero, in which case PE~0 falls idle early and its counter
substantially understates end-to-end latency. This mistake corrupted an
earlier measurement campaign and is documented here so it is not repeated.

% =====================================================================
\section{Verification}
% =====================================================================

Four layers. Each emitted testbench embeds an \sv{EXPECTED} constant computed
by the software oracle and asserts bit-equality of the observed root against
it, so every simulation is a self-checking test. Dynamic SystemVerilog
assertions cover the fabric invariants. A mutation harness injects two
classes of fault --- operand-address substitution and leaf corruption --- and
confirms the oracle flags them, establishing that the pass criterion is not
vacuous. Testbench-level SHA-256 hashing deduplicates functionally identical
images so that simulation effort is spent once per distinct program.

One subtlety in mutation testing deserves recording. A single-bit flip in a
leaf value is frequently \emph{not} observable at the root: at the operand
magnitudes used, a mantissa LSB is far below the round bit of the running sum
and round-to-nearest-even discards it. A harness that counts such mutants as
missed detections will report a misleadingly poor oracle. The correct
procedure is to escalate the perturbation until the software oracle confirms
the canonical root has actually changed, and only then require the hardware
oracle to flag it.

% =====================================================================
\section{Reproduction}
% =====================================================================

The compiler, RTL, benchmark suite, result CSVs and analysis scripts are
released with the main paper. The fabric elaborates under Vivado xsim 2025.2;
no vendor IP is instantiated, so it should port to other simulators with
minor adjustment, though only xsim has been exercised.

Two practical notes for anyone rerunning the campaign. Route-table
initialisation via the clocked \sv{write\_route()} protocol costs
$\mathcal{O}(G^2)$ simulated cycles and dominates wall-clock at large $G$;
rewriting those calls as hierarchical assignments is bit-exact for the root
result but not for cycle counts, and the substitution should be disclosed
wherever cycle figures are reported. And elaboration memory grows sharply
with $G$: the largest instance in the suite requires 50--70~GB per
elaboration, which is what put it out of reach of the campaign.

\end{document}
